\scanpage{105}
12) \uline{Relation entre les morphismes de générisation et morphismes de normalisation}

\deleted{Avant d'aborder une réalisation \unsure{calculatoire} de ce programme, j'ai envie d'examiner les relations entre les deux types de morphismes sur les $\Pi_1M_G$, $S\Pi_1M_G$ — \deleted{\unsure{les cas où les}}\inserted{\unsure{au delà de}} de la constatation $\pm$ triviale que pour des isomorphismes de graphes, les deux sont ``compatibles'' avec les projections $S\Pi_1M_G\longrightarrow\Pi_1M_G$. Pour le dire autrement, j'ai envie, pour une inclusion \uline{stricte} $I'\subset I$, d'examiner ``l'effet à l'infini'' des morphismes de groupoïdes (148).

Soit donc $G_{g,I}$ le \inserted{MD-}graphe réduit à un seul sommet de genre $g$, et avec $I$ comme ens de pts marqués. Soit $G$ un graphe de type numérique $(g,\nu=\operatorname{card} I)$, et considérons un morphisme de générisation
\[
(149)\qquad G\longrightarrow G_{g,I}\qquad\text{i.e. un iso. }I(G)\simeq I
\]
d'où un \deleted{diag.}\unsure{germe|gramme} de multiplicités
\margin{\deleted{$\unsure{M_G^*}$}}
\[
(150)\quad
\diagram{\begin{tikzcd}[ampersand replacement=\&, column sep=large,row sep=small]
M^*_{G,G_{g,I}} \arrow[r,leftrightarrow] \arrow[dr] \& (M_{G,G_{g,I}},\Gamma'_G) \& \\
\& M_{g,I} \& (M_G,\Gamma'_G)\arrow[ul,"\simeq"']
\end{tikzcd}}
\]
}
\margin{104}
